% easychair.tex,v 3.5 2017/03/15

\documentclass{easychair}

\usepackage{doc}

% use this if you have a long article and want to create an index
% \usepackage{makeidx}

% In order to save space or manage large tables or figures in a
% landcape-like text, you can use the rotating and pdflscape
% packages. Uncomment the desired from the below.
%
% \usepackage{rotating}
% \usepackage{pdflscape}

\makeatletter
\let\code\relax
\let\endcode\relax
\makeatother

\usepackage{agda}
\usepackage{newunicodechar}
\newunicodechar{⇒}{\ensuremath{\Rightarrow}}
\newunicodechar{◃}{\ensuremath{\triangleleft}}
\newunicodechar{▹}{\ensuremath{\triangleright}}
\newunicodechar{₁}{\ensuremath{_1}}
\newunicodechar{∈}{\ensuremath{\in}}
\newunicodechar{⊤}{\ensuremath{\top}}
\newunicodechar{≅}{\ensuremath{\cong}}

%\makeindex

%% Front Matter
%%
% Regular title as in the article class.
%
\title{Containers in Higher Kinds}

% Authors are joined by \and. Their affiliations are given by \inst, which indexes
% into the list defined using \institute
%
\author{
  Thorsten Altenkirch\inst{1} \and
  H{\aa}kon Robbestad Gylterud\inst{2} \and
  Zhili Tian\inst{1}
}

% Institutes for affiliations are also joined by \and,
\institute{
  $^{1}$School of Computer Science, University of Nottingham\\
  \email{\{psztxa,psxzt8\}@nottingham.ac.uk}\\
  $^{2}$Department of Informatics, University of Bergen\\
  \email{hakon.gylterud@uib.no}
}

%  \authorrunning{} has to be set for the shorter version of the authors' names;
% otherwise a warning will be rendered in the running heads. When processed by
% EasyChair, this command is mandatory: a document without \authorrunning
% will be rejected by EasyChair

\authorrunning{Altenkirch, Gylterud, Tian}

% \titlerunning{} has to be set to either the main title or its shorter
% version for the running heads. When processed by
% EasyChair, this command is mandatory: a document without \titlerunning
% will be rejected by EasyChair
\titlerunning{Containers in Higher Kinds}

\begin{document}

\maketitle

% \begin{abstract}
%   This is the abstract
% \end{abstract}

% The table of contents below is added for your convenience. Please do not use
% the table of contents if you are preparing your paper for publication in the
% EPiC Series or Kalpa Publications series

% \setcounter{tocdepth}{2}
% {\small
% \tableofcontents}

%\section{To mention}
%
%Processing in EasyChair - number of pages.
%
%Examples of how EasyChair processes papers. Caveats (replacement of EC
%class, errors).

\begin{code}[hide]%
\>[0]\AgdaSymbol{\{-\#}\AgdaSpace{}%
\AgdaKeyword{OPTIONS}\AgdaSpace{}%
\AgdaPragma{--guardedness}\AgdaSpace{}%
\AgdaPragma{--allow-unsolved-metas}\AgdaSpace{}%
\AgdaSymbol{\#-\}}\<%
\\
\>[0]\AgdaKeyword{module}\AgdaSpace{}%
\AgdaModule{main}\AgdaSpace{}%
\AgdaKeyword{where}\<%
\end{code}

\paragraph{Containers} 
Strictly positive types can be represented as containers \cite{containers},
that is, a shape \texttt{S : Set} and a family of positions \texttt{P : S → Set} giving rise to a
functor \texttt{⟦ S ◃ P ⟧ : Set ⇒ Set}. Every container has an initial algebra
(The W-type) and a terminal coalgebra (The M-type). However, there are
inductive/coinductive types which do not fit the scheme of W-types/M-types.
An example of such a type is the inductive type \texttt{Bush}, defined as follows:

\begin{code}%
\>[0]\AgdaKeyword{data}\AgdaSpace{}%
\AgdaDatatype{Bush}\AgdaSpace{}%
\AgdaSymbol{(}\AgdaBound{A}\AgdaSpace{}%
\AgdaSymbol{:}\AgdaSpace{}%
\AgdaPrimitive{Set}\AgdaSymbol{)}\AgdaSpace{}%
\AgdaSymbol{:}\AgdaSpace{}%
\AgdaPrimitive{Set}\AgdaSpace{}%
\AgdaKeyword{where}\<%
\\
\>[0][@{}l@{\AgdaIndent{0}}]%
\>[2]\AgdaInductiveConstructor{[]}\AgdaSpace{}%
\AgdaSymbol{:}\AgdaSpace{}%
\AgdaDatatype{Bush}\AgdaSpace{}%
\AgdaBound{A}\<%
\\
%
\>[2]\AgdaOperator{\AgdaInductiveConstructor{\AgdaUnderscore{}::\AgdaUnderscore{}}}\AgdaSpace{}%
\AgdaSymbol{:}\AgdaSpace{}%
\AgdaBound{A}\AgdaSpace{}%
\AgdaSymbol{→}\AgdaSpace{}%
\AgdaDatatype{Bush}\AgdaSpace{}%
\AgdaSymbol{(}\AgdaDatatype{Bush}\AgdaSpace{}%
\AgdaBound{A}\AgdaSymbol{)}\AgdaSpace{}%
\AgdaSymbol{→}\AgdaSpace{}%
\AgdaDatatype{Bush}\AgdaSpace{}%
\AgdaBound{A}\<%
\end{code}

Even though \texttt{Bush A} is not a W-type, the type constructor \texttt{Bush} itself arises
as the initial algebra of a functor on a functor category \cite{bird1998nested,
bird1999generalised}. This motivates the development of a notion of containers
at higher kinds (such as \texttt{(* ⇒ *) ⇒ (* ⇒ *)} or \texttt{(* ⇒ *) ⇒ *}) capable of modelling
strictly positive functors over containers.


\paragraph{Second-Order Containers}
We introduce a recursive syntax \texttt{2Cont} by extending normal containers with a family
\texttt{PF : S → Set} of positions of \texttt{F} and an inductive continuation \texttt{RF : (s : S) → PF s → 2Cont}.
These, second-order containers are of kind \texttt{(* ⇒ *) ⇒ (* ⇒ *)}. Hence, each of them gives rise to an endofunctor \texttt{⟦ S ◃ PX ◃ PF ◃ RF ⟧ : Cont ⇒ Cont}.

The \texttt{Bush} example can be specified using the following higher-order signature:
\texttt{H F X = 1 + X × F (F X)}. There is a sum involved and therefore two shapes
(\texttt{S = Bool}). The left shape is trivial with a constant, while in the right
shape, there is one position of \texttt{X} (\texttt{PX true = ⊤}) and one position of
\texttt{F} (\texttt{PF true = ⊤}). Finally, we need to model another second-order container
\texttt{H' F X = F X} which proceeds recursively (\texttt{RF true tt = H'}).

We define a second-order least fixpoint operator \texttt{2W : 2Cont → Cont} by
induction-recursion, thereby recovering \texttt{Bush} as the initial algebra of
\texttt{H}. The second-order greatest fixpoint perator \texttt{2M : 2Cont → Cont} 
can be derived similarly but using coinduction-induction.

\paragraph{Higher-Order Containers}
To continue the story of constructing strictly positive functors of functor
categories, we define a notion of higher-order container \texttt{HCont : Ty → Set}, where
\texttt{Ty} are just the types of simply typed $\lambda$-calculus with one base type
\texttt{*}, which represents \texttt{Set}. \texttt{HCont} is just a special case of
\texttt{Nf : Con → Ty → Set} where \texttt{Con} are the contexts of simply typed
$\lambda$-calculus: \texttt{HCont A = Nf • A}. We also use \texttt{Var : Con → Ty → Set}
for the typed de Bruijn variables. For brevity, we do not present the full
syntax. We remark that normal forms \texttt{Nf} are defined mutually with neutral terms
and spines, following the standard presentation of hereditary substitution for
simply typed $\lambda$-calculus \cite{keller2010normalization}. With this setting,
it is easy to see that \texttt{Set ≅ HCont *}, \texttt{Cont ≅ HCont (* ⇒ *)} and
\texttt{2Cont ≅ HCont ((* ⇒ *) ⇒ * ⇒ *)}. 

We can interpret every higher-order container as a function on types, \texttt{⟦\_⟧ :
HCont A → ⟦ A ⟧T} where \texttt{⟦\_⟧T  : Ty → Set₁} is the obvious
interpretation of types (with \texttt{⟦ * ⟧T = Set}). We believe that this can
be extended to hereditary functors.

However, using this semantics
we can show that there is no third-order fixed-point operator, \texttt{3W}.
Consider the following higher-order container\footnote{This is the corresponding 
$\lambda$-term, the actual encoding is not very readable.}:

\begin{code}[hide]%
\>[0]\AgdaKeyword{infixr}\AgdaSpace{}%
\AgdaNumber{20}\AgdaSpace{}%
\AgdaOperator{\AgdaInductiveConstructor{\AgdaUnderscore{}⇒\AgdaUnderscore{}}}\<%
\\
\>[0]\AgdaKeyword{data}\AgdaSpace{}%
\AgdaDatatype{Ty}\AgdaSpace{}%
\AgdaSymbol{:}\AgdaSpace{}%
\AgdaPrimitive{Set}\AgdaSpace{}%
\AgdaKeyword{where}\<%
\\
\>[0][@{}l@{\AgdaIndent{0}}]%
\>[2]\AgdaInductiveConstructor{*}\AgdaSpace{}%
\AgdaSymbol{:}\AgdaSpace{}%
\AgdaDatatype{Ty}\<%
\\
%
\>[2]\AgdaOperator{\AgdaInductiveConstructor{\AgdaUnderscore{}⇒\AgdaUnderscore{}}}\AgdaSpace{}%
\AgdaSymbol{:}\AgdaSpace{}%
\AgdaDatatype{Ty}\AgdaSpace{}%
\AgdaSymbol{→}\AgdaSpace{}%
\AgdaDatatype{Ty}\AgdaSpace{}%
\AgdaSymbol{→}\AgdaSpace{}%
\AgdaDatatype{Ty}\<%
\\
%
\\[\AgdaEmptyExtraSkip]%
\>[0]\AgdaKeyword{data}\AgdaSpace{}%
\AgdaDatatype{HCont}\AgdaSpace{}%
\AgdaSymbol{:}\AgdaSpace{}%
\AgdaDatatype{Ty}\AgdaSpace{}%
\AgdaSymbol{→}\AgdaSpace{}%
\AgdaPrimitive{Set}\AgdaSpace{}%
\AgdaKeyword{where}\<%
\end{code}

\begin{code}%
\>[0]\AgdaFunction{C}\AgdaSpace{}%
\AgdaSymbol{:}\AgdaSpace{}%
\AgdaDatatype{HCont}\AgdaSpace{}%
\AgdaSymbol{(((}\AgdaInductiveConstructor{*}\AgdaSpace{}%
\AgdaOperator{\AgdaInductiveConstructor{⇒}}\AgdaSpace{}%
\AgdaInductiveConstructor{*}\AgdaSymbol{)}\AgdaSpace{}%
\AgdaOperator{\AgdaInductiveConstructor{⇒}}\AgdaSpace{}%
\AgdaInductiveConstructor{*}\AgdaSymbol{)}\AgdaSpace{}%
\AgdaOperator{\AgdaInductiveConstructor{⇒}}\AgdaSpace{}%
\AgdaSymbol{(}\AgdaInductiveConstructor{*}\AgdaSpace{}%
\AgdaOperator{\AgdaInductiveConstructor{⇒}}\AgdaSpace{}%
\AgdaInductiveConstructor{*}\AgdaSymbol{)}\AgdaSpace{}%
\AgdaOperator{\AgdaInductiveConstructor{⇒}}\AgdaSpace{}%
\AgdaInductiveConstructor{*}\AgdaSymbol{)}\<%
\\
\>[0]\AgdaFunction{CC}\AgdaSpace{}%
\AgdaBound{F}\AgdaSpace{}%
\AgdaBound{G}\AgdaSpace{}%
\AgdaSymbol{=}\AgdaSpace{}%
\AgdaBound{G}\AgdaSpace{}%
\AgdaSymbol{(}\AgdaBound{F}\AgdaSpace{}%
\AgdaBound{G}\AgdaSymbol{)}\<%
\end{code}

\noindent
If an external fixpoint operator \texttt{3W : HCont (((* ⇒ *) ⇒ *) ⇒ (* ⇒ *) ⇒ *) → HCont ((* ⇒ *) ⇒ *)}
exists, applying it to \texttt{C} leads to an internal fixpoint operator \texttt{intW : HCont ((* ⇒ *) ⇒ *)}.
According to the semantics, \texttt{intW : (Set ⇒ Set) ⇒ Set} gives a least fixpoint of any functor, but we know such a thing does not
exist.

We would like to incorporate initial algebras and terminal coalgebras
at any type either by adding explicit fixpoint operators or by
allowing infinite terms (i.e. using a coinductively defined syntax).
In either case we will have to give up the naive semantics sketched
above. 

\paragraph{Higher-Order Containers as normalized $\lambda$-terms with $\Pi$ and $\Sigma$} 
Furthermore, we show higher-order containers are normalized $\lambda$-terms
closed under arbitrary products and coproducts (maybe also under least fixpoint and
greatest fixpoint). That is, we define a syntax:

\begin{code}[hide]%
\>[0]\AgdaKeyword{infixl}\AgdaSpace{}%
\AgdaNumber{5}\AgdaSpace{}%
\AgdaOperator{\AgdaInductiveConstructor{\AgdaUnderscore{}▹\AgdaUnderscore{}}}\<%
\\
\>[0]\AgdaKeyword{data}\AgdaSpace{}%
\AgdaDatatype{Con}\AgdaSpace{}%
\AgdaSymbol{:}\AgdaSpace{}%
\AgdaPrimitive{Set}\AgdaSpace{}%
\AgdaKeyword{where}\<%
\\
\>[0][@{}l@{\AgdaIndent{0}}]%
\>[2]\AgdaInductiveConstructor{•}%
\>[6]\AgdaSymbol{:}\AgdaSpace{}%
\AgdaDatatype{Con}\<%
\\
%
\>[2]\AgdaOperator{\AgdaInductiveConstructor{\AgdaUnderscore{}▹\AgdaUnderscore{}}}\AgdaSpace{}%
\AgdaSymbol{:}\AgdaSpace{}%
\AgdaDatatype{Con}\AgdaSpace{}%
\AgdaSymbol{→}\AgdaSpace{}%
\AgdaDatatype{Ty}\AgdaSpace{}%
\AgdaSymbol{→}\AgdaSpace{}%
\AgdaDatatype{Con}\<%
\\
%
\\[\AgdaEmptyExtraSkip]%
\>[0]\AgdaKeyword{data}\AgdaSpace{}%
\AgdaDatatype{Var}\AgdaSpace{}%
\AgdaSymbol{:}\AgdaSpace{}%
\AgdaDatatype{Con}\AgdaSpace{}%
\AgdaSymbol{→}\AgdaSpace{}%
\AgdaDatatype{Ty}\AgdaSpace{}%
\AgdaSymbol{→}\AgdaSpace{}%
\AgdaPrimitive{Set}\AgdaSpace{}%
\AgdaKeyword{where}\<%
\\
%
\\[\AgdaEmptyExtraSkip]%
\>[0]\AgdaKeyword{variable}\<%
\\
\>[0][@{}l@{\AgdaIndent{0}}]%
\>[2]\AgdaGeneralizable{Γ}\AgdaSpace{}%
\AgdaSymbol{:}\AgdaSpace{}%
\AgdaDatatype{Con}\<%
\\
%
\>[2]\AgdaGeneralizable{A}\AgdaSpace{}%
\AgdaGeneralizable{B}\AgdaSpace{}%
\AgdaSymbol{:}\AgdaSpace{}%
\AgdaDatatype{Ty}\<%
\end{code}

\begin{code}%
\>[0]\AgdaKeyword{data}\AgdaSpace{}%
\AgdaDatatype{Tm}\AgdaSpace{}%
\AgdaSymbol{:}\AgdaSpace{}%
\AgdaDatatype{Con}\AgdaSpace{}%
\AgdaSymbol{→}\AgdaSpace{}%
\AgdaDatatype{Ty}\AgdaSpace{}%
\AgdaSymbol{→}\AgdaSpace{}%
\AgdaPrimitive{Set₁}\AgdaSpace{}%
\AgdaKeyword{where}\<%
\\
\>[0][@{}l@{\AgdaIndent{0}}]%
\>[2]\AgdaInductiveConstructor{var}\AgdaSpace{}%
\AgdaSymbol{:}\AgdaSpace{}%
\AgdaDatatype{Var}\AgdaSpace{}%
\AgdaGeneralizable{Γ}\AgdaSpace{}%
\AgdaGeneralizable{A}\AgdaSpace{}%
\AgdaSymbol{→}\AgdaSpace{}%
\AgdaDatatype{Tm}\AgdaSpace{}%
\AgdaGeneralizable{Γ}\AgdaSpace{}%
\AgdaGeneralizable{A}\<%
\\
%
\>[2]\AgdaInductiveConstructor{lam}\AgdaSpace{}%
\AgdaSymbol{:}\AgdaSpace{}%
\AgdaDatatype{Tm}\AgdaSpace{}%
\AgdaSymbol{(}\AgdaGeneralizable{Γ}\AgdaSpace{}%
\AgdaOperator{\AgdaInductiveConstructor{▹}}\AgdaSpace{}%
\AgdaGeneralizable{A}\AgdaSymbol{)}\AgdaSpace{}%
\AgdaGeneralizable{B}\AgdaSpace{}%
\AgdaSymbol{→}\AgdaSpace{}%
\AgdaDatatype{Tm}\AgdaSpace{}%
\AgdaGeneralizable{Γ}\AgdaSpace{}%
\AgdaSymbol{(}\AgdaGeneralizable{A}\AgdaSpace{}%
\AgdaOperator{\AgdaInductiveConstructor{⇒}}\AgdaSpace{}%
\AgdaGeneralizable{B}\AgdaSymbol{)}\<%
\\
%
\>[2]\AgdaInductiveConstructor{app}\AgdaSpace{}%
\AgdaSymbol{:}\AgdaSpace{}%
\AgdaDatatype{Tm}\AgdaSpace{}%
\AgdaGeneralizable{Γ}\AgdaSpace{}%
\AgdaSymbol{(}\AgdaGeneralizable{A}\AgdaSpace{}%
\AgdaOperator{\AgdaInductiveConstructor{⇒}}\AgdaSpace{}%
\AgdaGeneralizable{B}\AgdaSymbol{)}\AgdaSpace{}%
\AgdaSymbol{→}\AgdaSpace{}%
\AgdaDatatype{Tm}\AgdaSpace{}%
\AgdaGeneralizable{Γ}\AgdaSpace{}%
\AgdaGeneralizable{A}\AgdaSpace{}%
\AgdaSymbol{→}\AgdaSpace{}%
\AgdaDatatype{Tm}\AgdaSpace{}%
\AgdaGeneralizable{Γ}\AgdaSpace{}%
\AgdaGeneralizable{B}\<%
\\
%
\>[2]\AgdaInductiveConstructor{Π}\AgdaSpace{}%
\AgdaSymbol{:}\AgdaSpace{}%
\AgdaSymbol{(}\AgdaBound{I}\AgdaSpace{}%
\AgdaSymbol{:}\AgdaSpace{}%
\AgdaPrimitive{Set}\AgdaSymbol{)}\AgdaSpace{}%
\AgdaSymbol{→}\AgdaSpace{}%
\AgdaSymbol{(}\AgdaBound{I}\AgdaSpace{}%
\AgdaSymbol{→}\AgdaSpace{}%
\AgdaDatatype{Tm}\AgdaSpace{}%
\AgdaGeneralizable{Γ}\AgdaSpace{}%
\AgdaGeneralizable{A}\AgdaSymbol{)}\AgdaSpace{}%
\AgdaSymbol{→}\AgdaSpace{}%
\AgdaDatatype{Tm}\AgdaSpace{}%
\AgdaGeneralizable{Γ}\AgdaSpace{}%
\AgdaGeneralizable{A}\<%
\\
%
\>[2]\AgdaInductiveConstructor{Σ}\AgdaSpace{}%
\AgdaSymbol{:}\AgdaSpace{}%
\AgdaSymbol{(}\AgdaBound{I}\AgdaSpace{}%
\AgdaSymbol{:}\AgdaSpace{}%
\AgdaPrimitive{Set}\AgdaSymbol{)}\AgdaSpace{}%
\AgdaSymbol{→}\AgdaSpace{}%
\AgdaSymbol{(}\AgdaBound{I}\AgdaSpace{}%
\AgdaSymbol{→}\AgdaSpace{}%
\AgdaDatatype{Tm}\AgdaSpace{}%
\AgdaGeneralizable{Γ}\AgdaSpace{}%
\AgdaGeneralizable{A}\AgdaSymbol{)}\AgdaSpace{}%
\AgdaSymbol{→}\AgdaSpace{}%
\AgdaDatatype{Tm}\AgdaSpace{}%
\AgdaGeneralizable{Γ}\AgdaSpace{}%
\AgdaGeneralizable{A}\<%
\end{code}

\noindent
together with a normalization function \texttt{nf : Tm Γ A → Nf Γ A}, 
which performs $\beta$-reduction and container normalization at the same time.
This mechanism allows one to derive higher-order container by spelling out its signature.

We also show \texttt{(Con, Tms, Ty, Tm)} is a model of simply-typed
$\lambda$-calculus (STLC), where \texttt{Tms} is a list of \texttt{Tm}. Similarly, we hope
to prove that the normalized model \texttt{(Con, Nfs, Ty, Nf)} is also a model of STLC, where
\texttt{Nfs} is a list of \texttt{Nf}.

\paragraph{Conclusion and Future Work}
Using Agda, we have formalized the syntax and semantics of second-order containers, including their least and greatest fixpoints. We have also derived a corresponding syntax for higher-order containers. However, this work remains ongoing, and several open questions persist.

In particular, we aim to demonstrate that the \texttt{W}-type remains sufficient. Specifically, for the second-order case, we conjecture that \texttt{2W} can be reduced to the first-order \texttt{W}, analogously to reducing \texttt{2M} to \texttt{M}. Additionally, we plan to develop a semantics for higher-order containers, together with their fixpoints and recursors.

\bibliographystyle{plain}
\bibliography{references}

%------------------------------------------------------------------------------
\end{document}